% Source-aligned English for printed p.125.
\SourcePage{125}

\LayerTarget{tohoku:section:1.2:p2}
\noindent transformation'' from $F$ to $G$ by certain authors) is a
``function'' that associates to every $A\in C$ a morphism $f(A)$ from $F(A)$
to $G(A)$, in such a way that for every morphism $u:A\longrightarrow B$ in
$C$, the following diagram
\[
\begin{array}{ccc}
F(A) & \stackrel{F(u)}{\longrightarrow} & F(B) \\
{\scriptstyle f(A)}\,\big\downarrow && \big\downarrow\,{\scriptstyle f(B)} \\
G(A) & \stackrel{G(u)}{\longrightarrow} & G(B)
\end{array}
\]
is commutative. Morphisms of functors $F\longrightarrow G$ and
$G\longrightarrow H$ compose in the evident way; this composition is
associative, and the ``identity morphism'' of the functor $F$ is an identity
for composition of morphisms of functors. (Thus, if $C$ is a set, the functors
from $C$ to $C'$ again form a category.) Finally, note that the composite $GF$
of two functors $F:C\longrightarrow C'$ and $G:C'\longrightarrow C''$ behaves
formally as a bifunctor with respect to the arguments $G$ and $F$: a functorial
morphism $G\longrightarrow G'$ (respectively, $F\longrightarrow F'$) defines a
functorial morphism $GF\longrightarrow G'F$ (respectively,
$GF\longrightarrow GF'$).

\LayerTarget{tohoku:section:1.2:p3}
An \textit{equivalence} of a category $C$ with a category $C'$ is a system
$(F,G,\varphi,\psi)$ formed of covariant functors
\[
  F:C\longrightarrow C' \qquad G:C'\longrightarrow C
\]
and homomorphisms of functors
\[
  \varphi:1_C\longrightarrow GF \qquad
  \psi:1_{C'}\longrightarrow FG,
\]
where $1_C$ and $1_{C'}$ are the identity functors in $C$ and $C'$,
respectively, such that for every $A\in C$ and $A'\in C'$ the composites
\[
\begin{array}{ccccc}
F(A) & \stackrel{F(\varphi(A))}{\longrightarrow} & FGF(A) &
\stackrel{\psi^{-1}(F(A))}{\longrightarrow} & F(A) \\
G(A') & \stackrel{G(\psi(A'))}{\longrightarrow} & GFG(A') &
\stackrel{\varphi^{-1}(G(A'))}{\longrightarrow} & G(A')
\end{array}
\]
are the identities of $F(A)$ and $G(A')$, respectively. Then for every pair
$A,B$ of objects of $C$, the map $f\longrightarrow F(f)$ from
$\mathrm{Hom}(A,B)$ to $\mathrm{Hom}(F(A),F(B))$ is a bijection, whose inverse
is the map $g\longrightarrow G(g)$ from $\mathrm{Hom}(F(A),F(B))$ to
$\mathrm{Hom}(GF(A),GF(B))$, identified with $\mathrm{Hom}(A,B)$ via the
isomorphisms $\varphi(A):A\longrightarrow GF(A)$ and
$\varphi(B):B\longrightarrow GF(B)$. Equivalences between categories compose
like functors. Two categories are said to be \textit{equivalent} if there
exists an equivalence between them. In ordinary language one then often allows
oneself not to distinguish between one and the other. It is important,
however, to observe the difference between this notion and the much stricter
notion of isomorphism (which applies if one wants to compare categories that
are sets): let $C$ be a nonempty set, and for every pair of objects $A,B\in C$
let $\mathrm{Hom}(A,B)$ be a singleton; then $C$ becomes (under the unique
possible composition laws
$\mathrm{Hom}(A,B)\times\mathrm{Hom}(B,C)\longrightarrow\mathrm{Hom}(A,C)$) a
category, and two categories constructed by this procedure are always
equivalent, but they are isomorphic only if they are equipotent. None of the
equivalences of categories encountered in practice is an isomorphism.

\TohokuContinuation{126}
